\section{Diseño experimental}

\subsection{Variables}

\begin{itemize}
    \item \textbf{Tamaño de la matriz} ($N$): 500, 1000, 2000 y 4000. Se eligió un rango con saltos de factor 2, desde problemas chicos (donde el paralelismo es marginal) hasta problemas grandes (donde el paralelismo muestra todo su potencial).
    \item \textbf{Cantidad de hilos} ($T$): 1 (secuencial), 2, 4, 8 y 16. Se usa exactamente el mismo conjunto que en la rama de pthreads del caso de estudio 1, para que las curvas sean directamente comparables.
    \item \textbf{Variable medida}: tiempo de la multiplicación en segundos, con \texttt{clock\_gettime} y \texttt{CLOCK\_MONOTONIC\_RAW}.
    \item \textbf{Métricas calculadas}: \textit{speedup} contra el secuencial, tasa de fallo de caché (\textit{miss rate}), número de instrucciones ejecutadas y memoria máxima residente (\textit{max RSS}).
\end{itemize}

En total se ejecutan \textbf{200 mediciones de desempeño} por máquina (5 valores de $T$ × 4 valores de $N$ × 10 repeticiones), más las 5 corrida de Cachegrind y las 5 de \texttt{getrusage} del perfilado.

\subsection{Metodología: 10 corridas por configuración}

Para cada combinación de $T$ y $N$ se ejecutan \textbf{10 corridas independientes} y se calcula la mediana. Como se explicó en la sección 5, una sola medición no es confiable por el ruido del sistema, y la mediana es el estimador robusto frente a valores atípicos.

\subsection{Orden de los bucles}

El script que automatiza el experimento recorre las repeticiones por fuera y los tamaños por dentro, de forma que las 10 repeticiones de un mismo $N$ no queden consecutivas. El motivo ya se explicó: evita que los datos queden \textit{calientes} en la caché y que la medición salga artificialmente rápida.

\subsection{Correctitud: validación contra una implementación de referencia}

Medir tiempos sin verificar que el programa calcula lo correcto es un error metodológico grave. Por eso, antes de cada corrida de medición, el script \texttt{experiments/medir\_openmp.sh} ejecuta una \textbf{validación de correctitud}:

\begin{enumerate}
    \item Se genera con \texttt{experiments/verify\_ref} una matriz resultado de referencia con $N=10$ y semilla 42, usando una implementación ingenua de tres bucles anidados.
    \item Se ejecuta el binario de OpenMP con los mismos parámetros ($N=10$, valor máximo 9, semilla 42).
    \item Se comparan las dos matrices $C$ elemento a elemento.
    \item Si \textbf{no coinciden exactamente}, el script aborta con un mensaje de error y no se mide nada.
\end{enumerate}

El mecanismo es el mismo que se usó en los labs 1 y 2 del caso de estudio 1, y por eso la implementación de referencia (\texttt{experiments/verify\_ref.c}) se reutilizó sin modificaciones.

\subsection{Reproducibilidad del benchmark entre máquinas}

El benchmark entre equipos se automatizó con \texttt{experiments/bench.sh}, que es un script \emph{portable}: corre en cualquier equipo con Linux sin necesidad de nada instalado que no venga de base, y produce un archivo de resultados con el hardware detectado más las corridas crudas. La guía completa para ejecutarlo en el segundo equipo está en \texttt{experiments/REPRODUCIR\_OTRA\_MAQUINA.md}.

\subsection{Herramientas}

\begin{center}
\begin{tabular}{|l|p{8.5cm}|}
\hline
\textbf{Herramienta} & \textbf{Uso en este trabajo} \\
\hline
\texttt{gcc} 14.2.0 & Compilador. Flag \texttt{-Wall -O2 -fopenmp}. \\
\texttt{gprof} & Perfilado de CPU por muestreo con instrumentación \texttt{-pg}. \\
\texttt{valgrind}/\texttt{cachegrind} & Perfilado de memoria: simulación de caché, D1/LL misses. \\
\texttt{getrusage} & Memoria máxima residente (max RSS) del proceso hijo. \\
\texttt{lscpu}, \texttt{free}, \texttt{uname} & Captura automática del hardware de cada equipo. \\
\texttt{matplotlib} & Generación de las gráficas de los resultados. \\
\hline
\end{tabular}
\end{center}
